% !TeX root = ../数学分析教材（伍胜健）笔记.tex
% 本文件由主文档通过 \input 载入；请编译主文档。
\setcounter{chapter}{2} % 保留原书章号。
\chapter{函数的极限与连续性}
\label{chap:function-limits-and-continuity}

序列极限研究下标趋于无穷时各项的变化，函数极限则允许自变量从不同方向趋于有限点或无穷远。\emph{二者的共同核心是：通过限制自变量的变化范围，使函数值与目标值的误差任意小。}本章先建立函数极限及其判别方法，再用极限刻画连续性；随后借助实数系的连续性，把逐点的局部结论推广为闭区间上的整体结论，最后用无穷小量与无穷大量的阶描述不同变化速度。

\section{函数的极限}
\label{sec:function-limits}

\subsection{函数极限的定义}
\label{subsec:definition-of-function-limits}

函数在某点可能没有定义，但在该点附近仍可能表现出确定的变化趋势。例如，$x\sin(1/x)$ 在 $x=0$ 没有定义，而 $|x\sin(1/x)|\leq|x|$ 表明它的值随着 $x\to0$ 趋近于零。这里所要刻画的是附近的函数值，而不是中心点的取值。以下沿用定义~\ref{def:real-neighborhood} 中的邻域 $U(x_0,\delta)$ 与去心邻域 $U_0(x_0,\delta)$。

\begin{definition}
\label{def:finite-function-limit}
设实函数 $f$ 在某个 $U_0(x_0,\delta_0)$ 上有定义，其中 $\delta_0>0$。若存在实数 $A$，使得
\begin{equation}
    \forall\varepsilon>0,\quad\exists\delta\in(0,\delta_0],\quad
    0<|x-x_0|<\delta\ \Longrightarrow\ |f(x)-A|<\varepsilon,
    \label{eq:epsilon-delta-function-limit}
\end{equation}
则称 $x$ 趋于 $x_0$ 时 $f(x)$ 的\emph{极限}（limit）为 $A$，记为
\[
    \lim_{x\to x_0}f(x)=A,
    \qquad\text{或}\qquad f(x)\to A\quad(x\to x_0).
\]
没有特别说明时，“极限存在”指有限极限存在。
\end{definition}

式~\eqref{eq:epsilon-delta-function-limit} 中，$\varepsilon$ 是任意给定的函数值误差，$\delta$ 则是为达到这一误差而选取的自变量范围。$\delta$ 可以依赖 $\varepsilon$、$x_0$ 与函数，但不能在选定之后再依赖邻域中的具体点 $x$。\emph[caution]{条件 $0<|x-x_0|$ 排除了中心点，所以极限与 $f(x_0)$ 是否有定义、取何值无关。}从几何上看，图像在足够小的去心邻域内全部落入水平带 $A-\varepsilon<y<A+\varepsilon$，而不只是有若干点落入其中。

\begin{example}
\label{ex:epsilon-delta-function-limits}
用定义证明
\[
    \lim_{x\to1}\frac{\sqrt{x}-1}{x-1}=\frac12,
    \qquad
    \lim_{x\to0}x\sin\frac1x=0.
\]
\end{example}

\begin{solve}
第一式先把差化为含 $|x-1|$ 的表达式。当 $0<|x-1|<1$ 时，$x>0$，从而
\[
    \left|\frac{\sqrt{x}-1}{x-1}-\frac12\right|
    =\left|\frac1{\sqrt{x}+1}-\frac12\right|
    =\frac{|x-1|}{2(\sqrt{x}+1)^2}
    \leq\frac{|x-1|}{2}.
\]
给定 $\varepsilon>0$，取 $\delta=\min\{1,2\varepsilon\}$，便有所需误差小于 $\varepsilon$。第二式由
\[
    \left|x\sin\frac1x-0\right|\leq|x|
\]
可直接取 $\delta=\varepsilon$。两式都不需要在极限点补定义函数值。
\end{solve}

\emph[tip]{用定义证明极限时，通常先限制 $|x-x_0|<\delta_1$，控制其余因子的大小，再把 $|f(x)-A|$ 估计为 $C|x-x_0|$ 或另一容易控制的量，最后取若干限制的最小值。}例如，$|x-2|<1$ 时 $|x|<3$，故
\[
    |x^4-16|=|x-2|\,|x+2|\,(x^2+4)\leq65|x-2|,
\]
从而证明 $\lim_{x\to2}x^4=16$ 时可取 $\delta=\min\{1,\varepsilon/65\}$。

\subsection{函数极限的性质}
\label{subsec:properties-of-function-limits}

函数极限只由任意充分小的去心邻域内的取值决定：如果两个函数在某个去心邻域内相同，它们在中心点的极限情况就相同。与收敛序列整体有界不同，函数有有限极限只能推出相应的\emph{局部有界性}，不能推出整个定义域上的有界性。

\begin{theorem}[唯一性、局部有界性与保序性]
\label{thm:local-properties-of-function-limits}
设 $\lim_{x\to x_0}f(x)=A\in\mathbb{R}$，则该极限唯一，且 $f$ 在 $x_0$ 的某个去心邻域内有界。若另有 $\lim_{x\to x_0}g(x)=B\in\mathbb{R}$，则有以下结论：若某个去心邻域内恒有 $f(x)\leq g(x)$，则 $A\leq B$；若 $A<B$，则对每个 $C\in(A,B)$，存在去心邻域使其中恒有 $f(x)<C<g(x)$。特别地，$A>0$ 时充分靠近 $x_0$ 的函数值为正，$A<0$ 时为负。
\end{theorem}

\begin{proof}
若 $A,B$ 都是 $f$ 的极限，则对任意 $\varepsilon>0$，在足够小的去心邻域内同时有 $|f(x)-A|<\varepsilon/2$ 与 $|f(x)-B|<\varepsilon/2$，故 $|A-B|<\varepsilon$。由 $\varepsilon$ 的任意性得 $A=B$。在极限定义中取 $\varepsilon=1$，得到局部估计 $|f(x)|<|A|+1$。

若 $A<B$，取 $\varepsilon=\tfrac12\min\{C-A,B-C\}$，并使两个误差同时小于 $\varepsilon$，便有 $f(x)<C<g(x)$。如果 $f(x)\leq g(x)$ 在某个去心邻域内成立但 $A>B$，将这一结论用于 $g,f$，就会在更小的去心邻域内得到 $g(x)<f(x)$，矛盾。因此 $A\leq B$。
\end{proof}

\emph[caution]{严格不等号在取极限后可能变成等号。}例如，$x^2>0$ 对所有 $x\neq0$ 成立，但 $\lim_{x\to0}x^2=0$。此外，$1/x$ 在每个 $x_0>0$ 处都有有限极限，却在 $(0,+\infty)$ 上无界。

\begin{theorem}[四则运算与夹逼原理]
\label{thm:function-limit-arithmetic-and-squeeze}
设 $f(x)\to A$、$g(x)\to B$（$x\to x_0$），其中 $A,B\in\mathbb{R}$，则
\begin{equation}
    \begin{aligned}
        \lim_{x\to x_0}(f(x)\pm g(x))&=A\pm B,\\
        \lim_{x\to x_0}f(x)g(x)&=AB,\\
        \lim_{x\to x_0}\frac{f(x)}{g(x)}&=\frac AB\qquad(B\neq0).
    \end{aligned}
    \label{eq:function-limit-arithmetic}
\end{equation}
商的条件 $B\neq0$ 保证 $g(x)$ 在某个去心邻域内不为零。若 $f(x)\leq h(x)\leq g(x)$ 在某个去心邻域内成立，且 $A=B$，则 $h(x)\to A$。
\end{theorem}

\begin{proof}
和与差的结论由误差不超过 $|f(x)-A|+|g(x)-B|$ 得到。乘积利用 $f$ 的局部有界性 $|f(x)|\leq M$，将误差估计为
\[
    |f(x)g(x)-AB|\leq M|g(x)-B|+|B|\,|f(x)-A|,
\]
令右端两个误差分别足够小即可。若 $B\neq0$，先限制 $|g(x)-B|<|B|/2$，则 $|g(x)|>|B|/2$，因而
\[
    \left|\frac1{g(x)}-\frac1B\right|
    \leq\frac{2}{|B|^2}|g(x)-B|\to0.
\]
再用乘积法则得到商的结论。对于夹逼原理，给定 $\varepsilon>0$，使两个外侧函数同时满足 $A-\varepsilon<f(x)$ 与 $g(x)<A+\varepsilon$，便有 $|h(x)-A|<\varepsilon$。
\end{proof}

\begin{theorem}[复合函数的极限]
\label{thm:composition-of-function-limits}
设 $g(x)\to u_0$（$x\to x_0$），$f(u)\to A$（$u\to u_0$），且复合函数在 $x_0$ 的某个去心邻域内有定义。如果该邻域内还有 $g(x)\neq u_0$，则
\begin{equation}
    \lim_{x\to x_0}f(g(x))=A.
    \label{eq:composition-of-function-limits}
\end{equation}
若 $f$ 在 $u_0$ 处有定义且 $f(u_0)=A$，则可以取消 $g(x)\neq u_0$ 这一条件。
\end{theorem}

\begin{proof}
给定 $\varepsilon>0$，由外层函数的极限选取 $\eta>0$，使 $0<|u-u_0|<\eta$ 时 $|f(u)-A|<\varepsilon$。再由内层函数的极限选取 $\delta>0$，使 $0<|x-x_0|<\delta$ 时 $|g(x)-u_0|<\eta$。在 $g(x)\neq u_0$ 的条件下，可以将 $u=g(x)$ 代入前一估计。若 $f(u_0)=A$，则 $g(x)=u_0$ 时误差本来就是零，结论仍成立。
\end{proof}

\begin{example}
\label{ex:composition-limit-exception}
说明复合函数极限定理中排除 $g(x)=u_0$ 的条件不能无条件删去。
\end{example}

\begin{solve}
设
\[
    f(u)=\begin{cases}1,&u\neq0,\\0,&u=0,\end{cases}
    \qquad
    g(x)=\begin{cases}x,&x\in\mathbb{Q},\\0,&x\notin\mathbb{Q}.\end{cases}
\]
由 $|g(x)|\leq|x|$ 得 $g(x)\to0$，且 $\lim_{u\to0}f(u)=1$。然而在非零有理数点，$f(g(x))=1$；在无理数点，$f(g(x))=0$。有理数与无理数都在实数中稠密，所以复合函数在 $0$ 附近始终取这两个不同的值，其极限不存在。问题在于内层函数反复取到 $0$，而 $f(0)$ 与外层函数的极限不同。
\end{solve}

四则运算应在核对各部分极限及分母条件后使用。例如，对正整数 $n$，先在 $x\neq1$ 时作恒等变形
\[
    \frac{x^n-1}{x-1}=1+x+\cdots+x^{n-1},
\]
再取极限，得到 $\lim_{x\to1}(x^n-1)/(x-1)=n$；原式中分母趋于零，不能直接套用商的法则。

\subsection{函数极限概念的推广}
\label{subsec:generalizations-of-function-limits}

自变量除了从两侧趋于一个有限点，还可以只从一侧趋近，或趋于无穷远；函数值也可能不趋于有限数而在绝对值上无限增大。推广定义时，关键是明确自变量最终进入什么集合，以及函数值需要满足什么条件。

\begin{definition}
\label{def:one-sided-function-limits}
设 $f$ 在 $(x_0,x_0+\delta_0)$ 上有定义。若
\[
    \forall\varepsilon>0,\quad\exists\delta\in(0,\delta_0],\quad
    0<x-x_0<\delta\ \Longrightarrow\ |f(x)-A|<\varepsilon,
\]
则称 $A$ 为 $f$ 在 $x_0$ 的\emph{右极限}（right-hand limit），记为 $\lim_{x\to x_0+}f(x)=A$，也记为 $f(x_0+0)=A$。将条件换成 $0<x_0-x<\delta$，便得到\emph{左极限}（left-hand limit），记为 $\lim_{x\to x_0-}f(x)=A$ 或 $f(x_0-0)=A$。
\end{definition}

\begin{proposition}[双侧极限与单侧极限]
\label{prop:two-sided-and-one-sided-limits}
设 $f$ 在 $x_0$ 的某个去心邻域内有定义，则
\[
    \lim_{x\to x_0}f(x)=A
    \quad\Longleftrightarrow\quad
    f(x_0-0)=f(x_0+0)=A.
\]
\end{proposition}

\begin{proof}
双侧极限的误差估计分别限制在左右两侧便给出必要性。反之，给定 $\varepsilon>0$，分别选取左右两侧的半径 $\delta_-,\delta_+$，再取 $\delta=\min\{\delta_-,\delta_+\}$，便在整个去心邻域内得到同一误差估计。
\end{proof}

例如，取整函数 $\lfloor x\rfloor$ 在整数 $k$ 处的左右极限分别为 $k-1,k$，所以双侧极限不存在。

\begin{definition}
\label{def:function-limits-at-infinity}
设 $f$ 在某个右侧无界区间上有定义。若
\[
    \forall\varepsilon>0,\quad\exists X>0,\quad
    x>X\ \Longrightarrow\ |f(x)-A|<\varepsilon,
\]
则记 $\lim_{x\to+\infty}f(x)=A$。相应地，将 $x>X$ 换成 $x<-X$ 或 $|x|>X$，分别定义 $\lim_{x\to-\infty}f(x)=A$ 与 $\lim_{x\to\infty}f(x)=A$，并要求函数在对应的尾部有定义。这里不带正负号的 $x\to\infty$ 表示 $|x|\to+\infty$，同时考察正、负两个方向。
\end{definition}

因此，$\lim_{x\to\infty}f(x)=A$ 当且仅当 $x\to+\infty$ 与 $x\to-\infty$ 的两个极限均为 $A$。上述有限极限仍有唯一性、相应尾部上的有界性、保序性、四则运算与夹逼性质，只需在证明中把去心邻域换成对应的自变量范围。

\begin{definition}
\label{def:generalized-function-limits}
设 $f$ 在 $x_0$ 的某个去心邻域内有定义。若
\[
    \forall M>0,\quad\exists\delta>0,\quad
    0<|x-x_0|<\delta\ \Longrightarrow\ f(x)>M,
\]
则称 $f(x)$ 趋于正无穷，记为 $\lim_{x\to x_0}f(x)=+\infty$。将结论换成 $f(x)<-M$ 或 $|f(x)|>M$，分别定义 $f(x)\to-\infty$ 与 $f(x)\to\infty$。这些称为\emph{广义极限}（generalized limits），最后一种也称函数为该过程中的\emph{无穷大量}（infinitely large quantity）。自变量趋向单侧或无穷远时，按相同方式替换自变量条件。
\end{definition}

\emph[caution]{函数值无界不等于函数趋于无穷大。}无界只要求能找到任意大的绝对值，而趋于无穷大要求自变量充分接近目标以后，所有函数值的绝对值都超过给定的界。例如，$x^2\sin x$ 在 $x\to+\infty$ 时无界，但在 $x=n\pi$ 处恒为零，故不是无穷大量。另一个区别是：$1/x$ 在 $x\to0$ 时满足 $|1/x|\to+\infty$，却既不趋于 $+\infty$，也不趋于 $-\infty$；其左右广义极限分别为 $-\infty,+\infty$。无穷符号不是实数，$\infty-\infty$、$0\cdot\infty$、$0/0$、$\infty/\infty$ 等形式不能按实数四则运算处理。

\begin{example}
\label{ex:radical-limits-at-infinity}
设 $F(x)=\sqrt{4x^2+2x+3}$，求 $x\to\pm\infty$ 时 $F(x)/x$ 的极限，以及
\[
    \lim_{x\to+\infty}(F(x)-2x),
    \qquad \lim_{x\to-\infty}(F(x)+2x).
\]
\end{example}

\begin{solve}
由 $\sqrt{x^2}=|x|$，有
\[
    \frac{F(x)}x=\frac{|x|}x\sqrt{4+\frac2x+\frac3{x^2}},
\]
故正、负方向的极限分别为 $2,-2$。对两个差式作有理化，得到
\[
    F(x)-2x=\frac{2x+3}{F(x)+2x},
    \qquad
    F(x)+2x=\frac{2x+3}{F(x)-2x}.
\]
在对应方向上将分子、分母同除以 $x$，分别得到 $2/(2+2)=1/2$ 与 $2/(-2-2)=-1/2$。这里既要消去无穷大量相减造成的不定形式，也要保留负方向上 $|x|/x=-1$ 的符号。
\end{solve}

\subsection{序列极限与函数极限的关系}
\label{subsec:sequence-and-function-limits}

要严格说明 $f(x)$ 在 $x\to x_0$ 时不趋于一个给定的实数 $A$，应将式~\eqref{eq:epsilon-delta-function-limit} 的量词逐个否定：
\begin{equation}
    \exists\varepsilon_0>0,\quad\forall\delta>0,\quad
    \exists x\in U_0(x_0,\delta),\quad |f(x)-A|\geq\varepsilon_0,
    \label{eq:negation-of-function-limit}
\end{equation}
其中只需考虑足够小且位于定义范围内的邻域。这表示任意靠近 $x_0$ 都仍能找到误差较大的点，并不要求整个邻域内的点都具有较大误差；它也只否定以 $A$ 为极限，尚未排除其他候选极限。

\begin{theorem}[海涅定理]
\label{thm:heine-sequential-criterion}
设 $f$ 在 $U_0(x_0,\delta_0)$ 上有定义，$A\in\mathbb{R}$。则 $\lim_{x\to x_0}f(x)=A$ 的充分必要条件是：对每一个满足
\[
    x_n\in U_0(x_0,\delta_0),\qquad x_n\to x_0
\]
的序列，都有 $f(x_n)\to A$。这一结论称为\emph{海涅定理}（Heine's theorem），也称函数极限的\emph{序列判别法}（sequential criterion）。
\end{theorem}

\begin{proof}
若函数极限为 $A$，给定 $\varepsilon>0$，先由函数极限得到 $\delta>0$，再由 $x_n\to x_0$ 得到充分大的 $n$ 满足 $0<|x_n-x_0|<\delta$，于是 $|f(x_n)-A|<\varepsilon$。

反之，若函数不以 $A$ 为极限，由式~\eqref{eq:negation-of-function-limit}，存在固定的 $\varepsilon_0>0$，对每个正整数 $n$ 都可选出 $x_n$ 满足
\[
    0<|x_n-x_0|<\min\{\delta_0,1/n\},
    \qquad |f(x_n)-A|\geq\varepsilon_0.
\]
这样 $x_n\to x_0$，而 $f(x_n)$ 不趋于 $A$，与假设矛盾。
\end{proof}

定理同样适用于单侧极限、无穷远处的极限以及广义极限，只需要求测试序列遵守相应的趋近方向。\emph[caution]{必须检查所有允许的序列；只检查一个序列，即使是 $x_n=x_0+1/n$，也不足以确认函数极限。}反过来，要否定极限存在，只需找出两条允许的序列，使对应函数值趋于不同的数，或者找出一条使函数值不收敛的序列。

\begin{example}
\label{ex:oscillatory-function-no-limit}
证明 $\lim_{x\to0}\sin(1/x)$ 不存在。
\end{example}

\begin{solve}
取 $x_n=1/(n\pi)$ 与 $y_n=1/(2n\pi+\pi/2)$，则两者都从右侧趋于零，而
\[
    \sin\frac1{x_n}=0,\qquad \sin\frac1{y_n}=1.
\]
由定理~\ref{thm:heine-sequential-criterion}，右极限已不存在，因而双侧极限也不存在。这与 $x\sin(1/x)\to0$ 并不矛盾：后者的振幅被趋于零的因子 $x$ 控制住了。
\end{solve}

\subsection{极限存在性定理和两个重要极限}
\label{subsec:limit-existence-and-two-important-limits}

极限的定义适合验证候选值，存在性定理则可以在尚未知道极限值时先判断它是否存在。单调性用确界确定单侧极限，柯西条件则通过比较邻近点的函数值来避免预先指定极限。

\begin{theorem}[单调函数的单侧极限]
\label{thm:monotone-function-one-sided-limits}
设 $f$ 在 $(a,b)$ 上单调递增，其中 $a,b$ 可以是无穷端点，则在扩展实数意义下
\begin{equation}
    \lim_{x\to a+}f(x)=\inf_{a<x<b}f(x),
    \qquad
    \lim_{x\to b-}f(x)=\sup_{a<x<b}f(x).
    \label{eq:monotone-function-endpoint-limits}
\end{equation}
无穷端点处分别理解为 $x\to-\infty$ 或 $x\to+\infty$。单调递减时，两个公式中的下确界、上确界互换。特别地，实值单调函数在区间每个内点都有有限的左右极限。
\end{theorem}

\begin{proof}
以递增函数的左端极限为例。记 $L=\inf f((a,b))$。若 $L$ 有限，由下确界性质，对任意 $\varepsilon>0$，存在 $c\in(a,b)$ 使 $L\leq f(c)<L+\varepsilon$。当 $a<x<c$ 时，由单调性有 $L\leq f(x)\leq f(c)<L+\varepsilon$，故左端极限为 $L$。若 $L=-\infty$，则对每个 $M>0$，可选 $c$ 使 $f(c)<-M$，于是 $a<x<c$ 时均有 $f(x)<-M$。右端极限及递减情形同理。在任一内点附近，用该点两旁固定位置的函数值夹住函数，便排除了无穷单侧极限。
\end{proof}

\begin{theorem}[函数极限的柯西收敛准则]
\label{thm:cauchy-criterion-for-function-limits}
设 $f$ 在 $x_0$ 的某个去心邻域内有定义，则 $\lim_{x\to x_0}f(x)$ 存在且有限，当且仅当
\begin{equation}
    \forall\varepsilon>0,\quad\exists\delta>0,\quad
    x,y\in U_0(x_0,\delta)\ \Longrightarrow\ |f(x)-f(y)|<\varepsilon.
    \label{eq:cauchy-condition-for-function-limits}
\end{equation}
\end{theorem}

\begin{proof}
若极限为 $A$，使 $|f(x)-A|$ 与 $|f(y)-A|$ 都小于 $\varepsilon/2$ 即得必要性。反之，固定一条位于定义范围内且趋于 $x_0$ 的序列 $\{x_n\}$，由式~\eqref{eq:cauchy-condition-for-function-limits}，$\{f(x_n)\}$ 是柯西列。由定理~\ref{thm:cauchy-convergence-criterion}，存在 $A\in\mathbb{R}$ 使 $f(x_n)\to A$。给定 $\varepsilon>0$，先将柯西条件中的误差取为 $\varepsilon/2$，得到 $\delta$；再固定足够大的 $n$，使 $x_n\in U_0(x_0,\delta)$ 且 $|f(x_n)-A|<\varepsilon/2$。于是对所有 $x\in U_0(x_0,\delta)$，
\[
    |f(x)-A|\leq|f(x)-f(x_n)|+|f(x_n)-A|<\varepsilon,
\]
所以 $f(x)\to A$。此证明只需选取一条固定序列，再利用柯西条件控制全部邻近点。
\end{proof}

\begin{theorem}[两个重要极限]
\label{thm:two-fundamental-function-limits}
当角度采用弧度制时，有
\begin{equation}
    \lim_{x\to0}\frac{\sin x}{x}=1,
    \qquad
    \lim_{t\to\infty}\left(1+\frac1t\right)^t=e.
    \label{eq:two-fundamental-function-limits}
\end{equation}
第二式在 $t\to+\infty$ 和 $t\to-\infty$ 两个方向均成立，且在对应尾部底数为正；等价地，
\begin{equation}
    \lim_{x\to0}(1+x)^{1/x}=e,
    \label{eq:fundamental-exponential-limit}
\end{equation}
其中取 $x>-1$ 且 $x\neq0$。
\end{theorem}

\begin{proof}
由单位圆中的三角形、扇形与外切三角形的面积比较，$0<x<\pi/2$ 时有 $\sin x<x<\tan x$，因此
\[
    \cos x<\frac{\sin x}{x}<1.
\]
这一比较也给出 $|\sin x|\leq|x|$；结合 $1-\cos x=2\sin^2(x/2)\leq x^2/2$，得到 $\cos x\to1$。由夹逼原理及 $\sin x/x$ 的偶性，第一个极限成立。弧度制使扇形面积恰为 $x/2$，所以该单位约定不能省略。

第二式以序列极限 $e=\lim_{n\to\infty}(1+1/n)^n$ 为基础。这个序列的收敛性也可由二项式展开核实：其第 $k$ 个展开项为 $\frac1{k!}\prod_{j=0}^{k-1}(1-j/n)$，随 $n$ 增大而不减，且由 $k!\geq2^{k-1}$（$k\geq1$），各项之和不超过 $\sum_{k=0}^n1/k!<3$，因而由单调收敛原理收敛。以下只用正底数实幂的代数性质与单调性，其构造与连续性见命题~\ref{prop:continuity-of-exponential-logarithmic-power-functions}。对实数 $t>1$，令 $n=\lfloor t\rfloor$，则
\[
    \left(1+\frac1{n+1}\right)^n
    \leq\left(1+\frac1t\right)^t
    \leq\left(1+\frac1n\right)^{n+1}.
\]
两端分别为 $\frac{(1+1/(n+1))^{n+1}}{1+1/(n+1)}$ 与 $(1+1/n)^n(1+1/n)$，均趋于 $e$，故正方向的极限为 $e$。当 $t\to-\infty$ 时，令 $s=-t-1\to+\infty$，则
\[
    \left(1+\frac1t\right)^t
    =\left(1+\frac1s\right)^{s+1}
    =\left(1+\frac1s\right)^s\left(1+\frac1s\right)\to e.
\]
最后令 $t=1/x$，便得到式~\eqref{eq:fundamental-exponential-limit}。
\end{proof}

\begin{example}
\label{ex:fundamental-trigonometric-limit-use}
求 $\lim_{x\to0}\sin(\sin x)/x$。
\end{example}

\begin{solve}
在充分小的去心邻域内 $\sin x\neq0$，并且 $\sin x\to0$，故
\[
    \frac{\sin(\sin x)}x
    =\frac{\sin(\sin x)}{\sin x}\frac{\sin x}x\to1\cdot1=1.
    \qedhere
\]
\end{solve}

普通极限不存在时，还可像序列一样描述邻近取值的上下边界。设 $f$ 在 $U_0(x_0,\delta_0)$ 内有定义，对 $0<\delta<\delta_0$，令
\[
    l(\delta)=\inf_{0<|x-x_0|<\delta}f(x),
    \qquad h(\delta)=\sup_{0<|x-x_0|<\delta}f(x).
\]
确界允许取扩展实数。随着 $\delta$ 缩小，$l(\delta)$ 不减，$h(\delta)$ 不增，因而可以定义\emph{函数的下极限与上极限}（limit inferior and limit superior）
\begin{equation}
    \liminf_{x\to x_0}f(x)=\lim_{\delta\to0+}l(\delta),
    \qquad
    \limsup_{x\to x_0}f(x)=\lim_{\delta\to0+}h(\delta).
    \label{eq:function-liminf-limsup}
\end{equation}
由 $l(\delta)\leq f(x)\leq h(\delta)$ 及确界定义可得：有限极限 $\lim_{x\to x_0}f(x)=A$ 存在，当且仅当上下极限都等于同一个\emph{有限数} $A$。具体地，极限存在时，任给 $\varepsilon>0$，充分小的邻域满足 $A-\varepsilon\leq l(\delta)\leq h(\delta)\leq A+\varepsilon$；反之，上下界同时趋于 $A$ 就夹住了全部邻近函数值。若上下极限共同等于 $+\infty$ 或 $-\infty$，得到的是相应的广义极限。例如，$\sin(1/x)$ 在 $0$ 的上下极限分别为 $-1,1$，而 $1/x^2$ 的上下极限均为 $+\infty$。

\section{函数的连续与间断}
\label{sec:continuity-and-discontinuity}

\subsection{函数的连续与间断}
\label{subsec:continuity-and-discontinuity}

% 教材第 103--109 页；分章 PDF 第 21--27 页。
极限只描述附近的变化，连续性还要求这种变化与该点的函数值衔接。讨论连续性时必须把定义域和端点包括在内；同一个表达式限制在不同定义域上时，允许的趋近方向也会改变。

\begin{definition}
\label{def:continuity-at-a-point-and-on-intervals}
设 $f:I\to\mathbb{R}$，其中 $I$ 为区间，$x_0\in I$。若
\begin{equation}
    \forall\varepsilon>0,\quad\exists\delta>0,\quad
    x\in I,\ |x-x_0|<\delta
    \ \Longrightarrow\ |f(x)-f(x_0)|<\varepsilon,
    \label{eq:epsilon-delta-continuity}
\end{equation}
则称 $f$ 在 $x_0$ \emph{连续}（continuous），否则称它在该点\emph{间断}（discontinuous）。当 $x_0$ 为内点时，这等价于 $\lim_{x\to x_0}f(x)=f(x_0)$。若只要求右侧或左侧的点满足上述条件，分别称为\emph{右连续}（right-continuous）或\emph{左连续}（left-continuous），相应条件是 $f(x_0+0)=f(x_0)$ 或 $f(x_0-0)=f(x_0)$。

若 $f$ 在区间 $I$ 的每一点都连续，则称 $f$ 在 $I$ 上连续，记为 $f\in C(I)$。特别地，$f\in C([a,b])$ 要求它在 $(a,b)$ 上连续、在 $a$ 处右连续且在 $b$ 处左连续。
\end{definition}

连续性也可以写成增量形式：令 $\Delta x=x-x_0$、$\Delta y=f(x_0+\Delta x)-f(x_0)$，则连续就是 $\Delta x\to0$ 时 $\Delta y\to0$，其中 $x_0+\Delta x$ 必须在定义域内。由海涅定理~\ref{thm:heine-sequential-criterion}，它还等价于：对任意 $\{x_n\}\subseteq I$，只要 $x_n\to x_0$，就有 $f(x_n)\to f(x_0)$。这里允许 $x_n=x_0$，因为这时函数值误差已经为零。

\begin{definition}
\label{def:classification-of-discontinuities}
设 $x_0$ 是 $f$ 的内点间断点。若左右极限都存在且有限，则称 $x_0$ 为\emph{第一类间断点}（discontinuity of the first kind）；其中，左右极限相等而不等于 $f(x_0)$ 时，称为\emph{可去间断点}（removable discontinuity），左右极限不等时，称为\emph{跳跃间断点}（jump discontinuity）。若至少一个单侧有限极限不存在，则称为\emph{第二类间断点}（discontinuity of the second kind）；其中包括单侧趋于无穷以及振荡而无极限的情形。
\end{definition}

可去间断点可以且只能通过把该点函数值改为共同的单侧极限来消除，跳跃与第二类间断不能靠改变一个点的值消除。若原函数在 $x_0$ 没有定义，则严格说来 $x_0$ 不属于其定义域；此时若有限极限 $A$ 存在，可补定义 $f(x_0)=A$，得到该点的\emph{连续延拓}（continuous extension）。类似地，$f\in C((a,b))$ 能连续延拓到 $[a,b]$，当且仅当两个端点的内侧极限均存在且有限，且端点值由这两个极限唯一确定。

\begin{example}
\label{ex:types-of-discontinuities}
讨论下列函数在 $0$ 处的连续性与间断类型：
\[
    f(x)=\begin{cases}x\sin(1/x),&x\neq0,\\1,&x=0,\end{cases}
    \qquad g(x)=\lfloor x\rfloor,
\]
\[
    h(x)=\begin{cases}e^{1/x},&x\neq0,\\0,&x=0,\end{cases}
    \qquad q(x)=\begin{cases}\sin(1/x),&x\neq0,\\0,&x=0.\end{cases}
\]
\end{example}

\begin{solve}
$f$ 的左右极限都为 $0$，而 $f(0)=1$，故 $0$ 是可去间断点，改取 $f(0)=0$ 即可连续。$g(0-0)=-1$、$g(0+0)=g(0)=0$，故 $g$ 在 $0$ 右连续，但该点是跳跃间断点。$h(0-0)=0$、$h(0+0)=+\infty$，故它在 $0$ 左连续，但该点属于第二类间断点。$q$ 的两侧极限均因振荡而不存在，故 $0$ 也是第二类间断点。后二者说明，第二类间断既可能伴随无界性，也可能发生在有界函数中；关于指数函数的连续性及尾部极限，见第~\ref{subsec:continuity-of-elementary-functions} 小节。
\end{solve}

\begin{example}
\label{ex:riemann-function-continuity}
在 $(0,1)$ 上定义\emph{黎曼函数}（Riemann function，也称 Thomae function）
\[
    R(x)=\begin{cases}
        1/q,&x=p/q\in\mathbb{Q},\ p,q\text{ 为互素正整数},\\
        0,&x\notin\mathbb{Q}.
    \end{cases}
\]
证明它在每个无理数点连续，在每个有理数点具有可去间断。
\end{example}

\begin{solve}
固定 $x_0\in(0,1)$。给定 $\varepsilon>0$，取正整数 $N$ 使 $1/N<\varepsilon$。区间内分母不超过 $N$ 的既约分数只有有限个，故可选 $\delta>0$，使 $U_0(x_0,\delta)\subseteq(0,1)$ 且排除这些有限个分数中所有异于 $x_0$ 的点。于是在该去心邻域内，有理数点的分母均大于 $N$，无理数点的值为零，因此
\[
    0\leq R(x)<1/N<\varepsilon.
\]
所以对每个 $x_0$ 都有 $\lim_{x\to x_0}R(x)=0$。无理数点满足 $R(x_0)=0$，故连续；有理数点满足 $R(x_0)>0$，故为可去间断点。这一例子表明，函数可以在稠密的点集上连续，同时在另一个稠密点集上间断。
\end{solve}

\begin{proposition}[单调函数的间断类型]
\label{prop:monotone-function-discontinuities}
实值单调函数在区间内部只可能有跳跃间断点。
\end{proposition}

\begin{proof}
以递增函数为例。由定理~\ref{thm:monotone-function-one-sided-limits}，内点 $x_0$ 的两个单侧极限均存在且有限，并由单调性有
\[
    f(x_0-0)\leq f(x_0)\leq f(x_0+0).
\]
若左右极限相等，则三者相等，函数在该点连续。因此只要它不连续，就必有 $f(x_0-0)<f(x_0+0)$，即发生跳跃。递减情形同理。
\end{proof}

\subsection{连续函数的性质}
\label{subsec:properties-of-continuous-functions}

% 教材第 109--111 页；分章 PDF 第 27--29 页。
连续函数在一点附近的性质可以直接由函数极限导出，但此时中心点也必须计入。连续性还允许把函数运算与极限过程交换，这是许多极限计算和后续微积分论证的基础。

\begin{proposition}[连续性的局部性质与运算]
\label{prop:local-properties-and-operations-of-continuity}
若 $f$ 在 $x_0$ 连续，则它在该点某个相对邻域内有界；若 $f(x_0)>0$，则对任意 $c\in(0,f(x_0))$，存在相对邻域使其中恒有 $f(x)>c$，负值情形类似。若 $f,g$ 都在 $x_0$ 连续，则 $f\pm g$、$fg$、$|f|$ 在该点连续；$g(x_0)\neq0$ 时，$f/g$ 也在该点连续。
\end{proposition}

\begin{proof}
局部有界性取式~\eqref{eq:epsilon-delta-continuity} 中的误差为 $1$ 即得。保号性取误差为 $f(x_0)-c$ 即得。四则运算由式~\eqref{eq:function-limit-arithmetic} 导出；绝对值的连续性则由反向三角不等式
\[
    \bigl||f(x)|-|f(x_0)|\bigr|\leq|f(x)-f(x_0)|
\]
得到。这些论证也适用于端点的单侧连续性。
\end{proof}

因此，连续函数的有限和、有限积仍连续，且
\[
    \max\{f,g\}=\frac{f+g+|f-g|}{2},
    \qquad \min\{f,g\}=\frac{f+g-|f-g|}{2}
\]
也连续。这里的运算次数有限；无限个连续函数的极限是否连续，需要后续函数列理论中的额外条件。

\begin{theorem}[复合函数的连续性]
\label{thm:composition-of-continuous-functions}
设 $g$ 在 $x_0$ 连续，$f$ 在 $u_0=g(x_0)$ 连续，且 $g$ 在 $x_0$ 附近的取值属于 $f$ 的定义域，则 $f\circ g$ 在 $x_0$ 连续。
\end{theorem}

\begin{proof}
给定 $\varepsilon>0$，先由 $f$ 在 $u_0$ 连续选取 $\eta>0$，使 $|u-u_0|<\eta$ 时 $|f(u)-f(u_0)|<\varepsilon$；再由 $g$ 在 $x_0$ 连续选取 $\delta>0$，使 $|x-x_0|<\delta$ 时 $|g(x)-g(x_0)|<\eta$。代入即得 $|f(g(x))-f(g(x_0))|<\varepsilon$。由于外层函数在中心点也满足这一估计，不需要排除 $g(x)=u_0$。
\end{proof}

\begin{theorem}[反函数的连续性]
\label{thm:continuity-of-inverse-functions}
设 $I$ 为非退化区间，$f:I\to\mathbb{R}$ 连续且严格单调，则 $f$ 的值域 $J=f(I)$ 是区间，反函数 $f^{-1}:J\to I$ 连续且具有与 $f$ 相同的严格单调方向。端点处的连续性按相对定义域理解。
\end{theorem}

\begin{proof}
严格单调性保证单射性，从而反函数在 $J$ 上存在，且具有相同的单调方向。$J$ 是区间将由下一节的介值定理~\ref{thm:intermediate-value-property} 得到；下面的连续性证明只需对 $y\in J$ 考察，不依赖尚未证明的介值性。

不妨设 $f$ 严格递增。固定 $y_0=f(x_0)$，先设 $x_0$ 是 $I$ 的内点。给定 $\varepsilon>0$，在 $I$ 内选 $x_-,x_+$ 使
\[
    x_0-\varepsilon<x_-<x_0<x_+<x_0+\varepsilon.
\]
令 $\delta=\min\{y_0-f(x_-),f(x_+)-y_0\}>0$。若 $y\in J$ 且 $|y-y_0|<\delta$，由严格单调性得
\[
    x_-<f^{-1}(y)<x_+,
\]
从而 $|f^{-1}(y)-x_0|<\varepsilon$。若 $x_0$ 是属于 $I$ 的左端点，只选 $x_+$ 并利用 $f^{-1}(y)\geq x_0$；右端点类似。递减情形可对 $-f$ 作同样论证。
\end{proof}

\subsection{初等函数的连续性}
\label{subsec:continuity-of-elementary-functions}

% 教材第 111--113 页；分章 PDF 第 29--31 页。
初等函数的连续性由基本初等函数逐步建立，再通过有限次四则运算与复合传递。多项式的连续性来自 $x^n$ 的极限；有理函数在分母非零处连续；三角函数则可先由和差公式得到
\[
    |\sin x-\sin y|\leq|x-y|,
    \qquad |\cos x-\cos y|\leq|x-y|,
\]
从而正弦、余弦处处连续，其余三角函数在相应分母非零处连续。证明这些结论不需要预先使用导数。

对于实指数，先说明无理指数的含义。给定 $a>1$，有理数 $r=p/q$ 的幂由正数开方定义，负指数通过倒数定义；对任意实数 $x$，可统一写成
\begin{equation}
    a^x=\sup\{a^r:r\in\mathbb{Q},\ r<x\}.
    \label{eq:real-exponent-via-supremum}
\end{equation}
这个集合非空且有上界，因为可以选整数 $m<x<n$，使 $a^m$ 属于集合，而 $a^n$ 是上界。该定义与有理指数一致：由伯努利不等式~\eqref{eq:bernoulli-inequality}，$a^{1/N}\to1$，因为若 $a^{1/N}\geq1+\eta$，则 $a\geq(1+\eta)^N\geq1+N\eta$，对固定 $\eta>0$ 和充分大 $N$ 不可能成立。对有理数 $x$，从下方用 $x-1/N$ 逼近即可验证式~\eqref{eq:real-exponent-via-supremum} 的确等于原来的 $a^x$。

\begin{proposition}[指数、对数与幂函数的连续性]
\label{prop:continuity-of-exponential-logarithmic-power-functions}
对 $a>0$，指数函数 $a^x$ 在 $\mathbb{R}$ 上连续；$a\neq1$ 时，其反函数 $\log_a x$ 在 $(0,+\infty)$ 上连续。对任意 $\alpha\in\mathbb{R}$，幂函数 $x^\alpha=e^{\alpha\ln x}$ 在 $(0,+\infty)$ 上连续。
\end{proposition}

\begin{proof}
先设 $a>1$。由确界定义及有理数稠密性，若 $x_1<x_2$，可选有理数 $r,s$ 满足 $x_1<r<s<x_2$，得到 $a^{x_1}\leq a^r<a^s\leq a^{x_2}$，故指数函数严格递增。固定 $x_0\in\mathbb{R}$，对正整数 $N$，令
\[
    r_N=\frac{\lfloor Nx_0\rfloor-1}{N},
    \qquad s_N=\frac{\lfloor Nx_0\rfloor+2}{N}.
\]
则 $r_N<x_0<s_N$、$s_N-r_N=3/N$。当 $x$ 足够靠近 $x_0$，使 $r_N<x<s_N$ 时，$a^x,a^{x_0}$ 均在 $[a^{r_N},a^{s_N}]$ 内，因而
\[
    |a^x-a^{x_0}|\leq a^{s_N}-a^{r_N}
    =a^{r_N}(a^{3/N}-1)\leq a^m(a^{3/N}-1),
\]
其中 $m>x_0$ 为固定整数。已知 $a^{1/N}\to1$，故先选充分大的 $N$，再选相应邻域，就能使此误差任意小，证明连续性。$a=1$ 时为常函数；$0<a<1$ 时定义 $a^x=1/(1/a)^x$，由商的连续性得到结论。

当 $a>1$ 时，$a^n\to+\infty$、$a^{-n}\to0$，结合单调性和介值性，指数函数的值域为 $(0,+\infty)$；$0<a<1$ 时同理。因此由反函数连续性定理，对数函数连续。幂函数的连续性再由指数、对数的复合得到。
\end{proof}

严格单调的三角函数分支的反函数也连续：例如，$\arcsin x$ 在 $[-1,1]$ 上连续，端点处按单侧理解；$\arccos x$ 在 $[-1,1]$ 上连续，$\arctan x$ 在 $\mathbb{R}$ 上连续。于是由有限次运算得到的\emph{初等函数在其定义域内连续}。\emph[caution]{“在定义域内”是结论的一部分：例如，$\ln x$ 不在 $0$ 有定义，$1/x$ 的定义域不含 $0$，不能通过直接代入这些点求有限极限。}

\begin{example}
\label{ex:variable-exponent-function-limit}
设 $u(x)\to a>0$、$v(x)\to b\in\mathbb{R}$（$x\to x_0$），证明 $u(x)^{v(x)}\to a^b$，并求 $\lim_{x\to0}(1-2x)^{1/x}$。
\end{example}

\begin{solve}
由 $a>0$ 及局部保号性，$u(x)$ 在充分小的去心邻域内为正，可写成
\[
    u(x)^{v(x)}=e^{v(x)\ln u(x)}\to e^{b\ln a}=a^b.
\]
第二问的指数趋于无穷，不能使用刚证明的有限指数结论。令 $t=-2x$，则
\[
    (1-2x)^{1/x}=\bigl((1+t)^{1/t}\bigr)^{-2}\to e^{-2},
\]
其中先用式~\eqref{eq:fundamental-exponential-limit}，再用正数上的幂函数连续性。这说明处理变指数极限时，应先判断指数是否有有限极限。
\end{solve}

\section{闭区间上连续函数的基本性质}
\label{sec:continuous-functions-on-closed-intervals}

% 教材第 113--121 页；分章 PDF 第 31--39 页。
逐点连续给出的邻域大小和局部界可以随点变化，而闭区间的有限覆盖性质与有界序列的子列性质，使这些局部信息能够推出整体结论。以下除另有说明外，均设 $a,b\in\mathbb{R}$ 且 $a<b$；闭性与有界性都是区间条件中不可随意删去的部分。

\begin{theorem}[有界性与最值定理]
\label{thm:boundedness-and-extreme-values}
若 $f\in C([a,b])$，则 $f$ 在 $[a,b]$ 上有界，并存在 $\xi,\eta\in[a,b]$，使
\begin{equation}
    f(\xi)=\min_{x\in[a,b]}f(x)=m,
    \qquad f(\eta)=\max_{x\in[a,b]}f(x)=M.
    \label{eq:extreme-values-on-closed-interval}
\end{equation}
\end{theorem}

\begin{proof}
先证有界性。由连续性的局部有界性，每个 $x\in[a,b]$ 都有邻域 $U(x,\delta_x)$，使 $f$ 在该邻域与 $[a,b]$ 的交集上有界。对这些邻域应用有限覆盖定理~\ref{thm:finite-covering}，取有限子覆盖，再取相应有限多个局部界的最大值，就得到整个区间上的界。这也完成了例~\ref{ex:continuous-function-bounded-by-finite-cover} 中提出的证明。

由有界性，$M=\sup\{f(x):x\in[a,b]\}$ 是实数。由上确界定义，对每个正整数 $n$ 可选 $x_n\in[a,b]$，使 $M-1/n<f(x_n)\leq M$，于是 $f(x_n)\to M$。由定理~\ref{thm:bolzano-weierstrass-sequence}，存在子列 $x_{n_k}\to\eta$；闭区间条件保证 $\eta\in[a,b]$。再由连续性，
\[
    f(\eta)=\lim_{k\to\infty}f(x_{n_k})=M.
\]
对 $-f$ 应用同一论证，便得到最小值。
\end{proof}

\emph[caution]{有界只保证确界有限，最值定理还保证确界实际取到。}例如，$1/x$ 在 $(0,1)$ 上连续却无界，$x$ 在 $(0,1)$ 上连续且有界却没有最大值与最小值；$x$ 在 $\mathbb{R}$ 上连续也不能推出有界性。这些例子分别说明开区间和无界区间不足以替代有界闭区间。

\begin{corollary}[严格正的连续函数有统一正下界]
\label{cor:positive-continuous-function-lower-bound}
若 $f\in C([a,b])$ 且 $f(x)>0$ 对每个 $x\in[a,b]$ 成立，则存在 $c>0$ 使 $f(x)\geq c$ 对整个区间成立。特别地，若连续函数在闭区间上处处非零，则 $|f|$ 有正下界，$1/f$ 连续且有界。
\end{corollary}

\begin{proof}
由最值定理，存在 $\xi\in[a,b]$ 使 $c=f(\xi)=\min f$。由于最小值在区间内取到，且函数处处为正，故 $c>0$。对 $|f|$ 应用这一结论即得后半部分。
\end{proof}

\begin{theorem}[介值定理]
\label{thm:intermediate-value-property}
设 $f$ 在区间 $I$ 上连续，$x_1,x_2\in I$。若实数 $c$ 介于 $f(x_1)$ 与 $f(x_2)$ 之间，则存在介于 $x_1,x_2$ 之间的点 $\xi\in I$ 使 $f(\xi)=c$；当 $c$ 严格介于两个函数值之间时，可以取 $\xi$ 严格介于两个点之间。这一性质称为\emph{介值性}（intermediate value property）。因此 $f(I)$ 是区间；特别地，若 $f\in C([a,b])$，则
\begin{equation}
    f([a,b])=[m,M],
    \label{eq:continuous-image-of-closed-interval}
\end{equation}
其中 $m,M$ 由式~\eqref{eq:extreme-values-on-closed-interval} 给出。
\end{theorem}

\begin{proof}
端点值相等或 $c$ 等于端点值时直接取相应点。对于严格介于的情况，通过交换两点或考虑 $-f$，只需证明 $u<v$ 且 $f(u)<c<f(v)$ 的情形。令
\[
    E=\{x\in[u,v]:f(x)\leq c\},\qquad \xi=\sup E.
\]
$E$ 非空且有上界。由 $u,v$ 处的连续性与严格不等式，$u$ 右侧有一段仍满足 $f(x)<c$，$v$ 左侧有一段满足 $f(x)>c$，故 $u<\xi<v$。由确界性质，存在 $t_n\in E$ 使 $t_n\to\xi$，因此连续性给出 $f(\xi)=\lim f(t_n)\leq c$。若 $f(\xi)<c$，则由连续性，$\xi$ 右侧足够靠近的点仍满足 $f(x)<c$，从而属于 $E$，与 $\xi$ 为上界矛盾。故 $f(\xi)=c$。

任取 $f(I)$ 中两个值，它们之间的值都由上述结论取到，所以 $f(I)$ 是区间。闭区间情形再结合最大值、最小值的存在性，就得到式~\eqref{eq:continuous-image-of-closed-interval}。
\end{proof}

\begin{corollary}[零点存在定理]
\label{cor:bolzano-zero-existence}
若 $f\in C([a,b])$ 且 $f(a)f(b)<0$，则存在 $\xi\in(a,b)$ 使 $f(\xi)=0$。此结论也称\emph{波尔查诺零点定理}（Bolzano's theorem），由介值定理取中间值 $0$ 得到。
\end{corollary}

\emph[caution]{端点异号保证至少有一个零点，不保证唯一。}例如，$x(x^2-1)$ 在 $[-2,2]$ 上满足端点异号，却有三个零点。若另有严格单调性，则同一个函数值至多取到一次，可以把存在性与唯一性结合起来。介值定理还表明，对连续函数的有限个取值 $f(x_1),\ldots,f(x_n)$，其算术平均介于这些值的最小值与最大值之间，因而也必是区间内某点的函数值。

\begin{proposition}[连续函数的单射性与严格单调性]
\label{prop:continuous-injective-function-strict-monotonicity}
设 $f$ 在非退化区间 $I$ 上连续，则它为单射的充分必要条件是它严格单调。因此，区间上的连续函数在其值域上存在反函数，当且仅当它严格单调。
\end{proposition}

\begin{proof}
严格单调必为单射。反之，若 $f$ 连续且为单射，对任意 $x_1<x_2<x_3$，$f(x_2)$ 必严格介于 $f(x_1),f(x_3)$ 之间。否则它比两端值都大或都小，在中间函数值与两端函数值之间选取一个适当的值，由介值定理，该值会在 $(x_1,x_2)$ 和 $(x_2,x_3)$ 内各取到一次，违反单射性。

固定 $u<v$。若 $f(u)<f(v)$，任取 $s<t$，把 $u,v,s,t$ 的不同点按大小排列。上述三点性质保证其中任意相邻两段的函数值变化方向一致，而 $f(u)<f(v)$ 又确定这一方向只能递增，故 $f(s)<f(t)$。若 $f(u)>f(v)$，同理得到严格递减。
\end{proof}

对于一般的单调函数，连续性的另一判据是其值域是否为区间：连续性由介值定理保证值域为区间；反之，单调函数若发生间断，有限的单侧极限之间必有跳跃，从而遗漏一段值中的点，值域就不可能是区间。端点处同样考察函数值与内侧极限之间的跳跃。

\begin{example}
\label{ex:continuous-zero-and-bisection}
证明方程 $x^4+2x-1=0$ 在 $(0,1)$ 内有唯一解，并说明如何用二分法给出带误差保证的近似解。
\end{example}

\begin{solve}
令 $F(x)=x^4+2x-1$。它在 $[0,1]$ 上连续，且 $F(0)=-1<0<F(1)=2$，故存在零点。若 $0\leq x<y\leq1$，则
\[
    F(y)-F(x)=(y-x)(y^3+y^2x+yx^2+x^3+2)>0,
\]
故 $F$ 严格递增，零点 $\xi$ 唯一。

取初始区间 $[a_0,b_0]=[0,1]$。每步令 $c_n=(a_n+b_n)/2$；若 $F(c_n)=0$，便得到精确解并终止；否则保留使端点函数值异号的半区间。零点始终在保留区间内，进行 $n$ 次二分后 $b_n-a_n=2^{-n}$，所以其中点 $c_n$ 满足
\[
    |c_n-\xi|\leq\frac{b_n-a_n}{2}=2^{-(n+1)}.
\]
给定误差容限 $\varepsilon>0$，选取 $2^{-(n+1)}\leq\varepsilon$ 即可停止并输出 $c_n$。例如，八次二分后得到区间 $[121/256,122/256]$，其中点 $243/512=0.474609375$ 的误差不超过 $1/512$。二分法的根存在性来自连续与异号条件，误差保证来自区间长度；唯一性则由另外证明的严格单调性提供。
\end{solve}

逐点连续仍允许不同位置需要不同的邻域半径。为了在整个区间内统一控制函数值的变化，需要加强量词要求。

\begin{definition}
\label{def:uniform-continuity}
设 $f:I\to\mathbb{R}$。若
\begin{equation}
    \forall\varepsilon>0,\quad\exists\delta>0,\quad
    \forall x,y\in I,\quad
    |x-y|<\delta\ \Longrightarrow\ |f(x)-f(y)|<\varepsilon,
    \label{eq:uniform-continuity}
\end{equation}
则称 $f$ 在 $I$ 上\emph{一致连续}（uniformly continuous）。
\end{definition}

逐点连续的次序是“对每个 $x_0$，对每个 $\varepsilon$，存在 $\delta$”，而一致连续要求“对每个 $\varepsilon$，存在同时适用于所有点对的 $\delta$”。因此一致连续蕴含连续，但反向一般不成立。否定式为
\begin{equation}
    \exists\varepsilon_0>0,\quad\forall\delta>0,\quad
    \exists x,y\in I,\quad |x-y|<\delta,
    \quad |f(x)-f(y)|\geq\varepsilon_0.
    \label{eq:negation-of-uniform-continuity}
\end{equation}

\begin{theorem}[一致连续的序列判据]
\label{thm:sequential-criterion-uniform-continuity}
函数 $f$ 在区间 $I$ 上一致连续，当且仅当对任意两列 $x_n,y_n\in I$，只要 $|x_n-y_n|\to0$，就有 $|f(x_n)-f(y_n)|\to0$。这里不要求 $x_n,y_n$ 各自收敛。
\end{theorem}

\begin{proof}
必要性直接将式~\eqref{eq:uniform-continuity} 用于充分大的 $n$。若一致连续不成立，在式~\eqref{eq:negation-of-uniform-continuity} 中依次取 $\delta=1/n$，得到 $|x_n-y_n|<1/n$，而函数值之差始终不小于固定的 $\varepsilon_0$，与序列条件矛盾。
\end{proof}

\begin{example}
\label{ex:uniform-continuity-examples}
讨论 $\sin x$ 在 $\mathbb{R}$ 上、$1/x$ 在 $(0,+\infty)$ 与 $[c,+\infty)$（$c>0$）上，以及 $\sqrt{x}$ 在 $[0,+\infty)$ 上的一致连续性。
\end{example}

\begin{solve}
由 $|\sin x-\sin y|\leq|x-y|$，对任意 $\varepsilon>0$ 取 $\delta=\varepsilon$，得正弦函数在整条实轴上一致连续。对 $1/x$，取 $x_n=1/n$、$y_n=1/(n+1)$，则两点距离趋于零，但
\[
    \left|\frac1{x_n}-\frac1{y_n}\right|=1,
\]
故它在 $(0,+\infty)$ 上不一致连续；在 $[c,+\infty)$ 上则有
\[
    \left|\frac1x-\frac1y\right|=\frac{|x-y|}{xy}\leq\frac{|x-y|}{c^2},
\]
取 $\delta=c^2\varepsilon$ 即可。最后，对 $x,y\geq0$，不妨设 $x\geq y$，则
\[
    (\sqrt{x}-\sqrt{y})^2\leq(\sqrt{x}-\sqrt{y})(\sqrt{x}+\sqrt{y})=x-y,
\]
所以 $|\sqrt{x}-\sqrt{y}|\leq\sqrt{|x-y|}$，取 $\delta=\varepsilon^2$ 即得一致连续。这也说明在无界区间上一致连续的函数可以无界。
\end{solve}

\begin{theorem}[康托尔定理]
\label{thm:heine-cantor-uniform-continuity}
若 $f\in C([a,b])$，则 $f$ 在 $[a,b]$ 上一致连续。这一结论称为\emph{康托尔定理}（Heine--Cantor theorem）。
\end{theorem}

\begin{proof}
反设不一致连续。由式~\eqref{eq:negation-of-uniform-continuity}，可选两列 $x_n,y_n\in[a,b]$，使 $|x_n-y_n|<1/n$ 且 $|f(x_n)-f(y_n)|\geq\varepsilon_0>0$。由有界序列的子列定理，可选子列 $x_{n_k}\to\xi\in[a,b]$。又因
\[
    |y_{n_k}-\xi|\leq|y_{n_k}-x_{n_k}|+|x_{n_k}-\xi|\to0,
\]
故 $y_{n_k}\to\xi$。在 $\xi$ 处的连续性给出 $f(x_{n_k}),f(y_{n_k})\to f(\xi)$，于是两者之差趋于零，与固定的正下界矛盾。
\end{proof}

\begin{proposition}[有穷开区间上的一致连续性]
\label{prop:uniform-continuity-and-endpoint-extension}
设 $a,b$ 为有限实数，$f\in C((a,b))$。则 $f$ 在 $(a,b)$ 上一致连续，当且仅当两个单侧极限 $f(a+0),f(b-0)$ 都存在且有限，也即 $f$ 可以连续延拓到 $[a,b]$。
\end{proposition}

\begin{proof}
若一致连续，给定 $\varepsilon>0$，取相应的 $\delta>0$。任意两个位于 $(a,a+\eta)$ 的点，在 $0<\eta<\min\{\delta,b-a\}$ 时距离小于 $\delta$，其函数值之差便小于 $\varepsilon$。这正是 $x\to a+$ 的柯西条件，故由单侧版本的定理~\ref{thm:cauchy-criterion-for-function-limits}，右极限存在且有限；右端点同理。反之，以这两个极限补上端点值，得到闭区间上的连续函数，由康托尔定理知其一致连续，限制回 $(a,b)$ 后结论仍成立。
\end{proof}

因此，有穷区间上的一致连续函数必有界：开区间情形可通过上述延拓及最值定理看出，其他端点类型同理。但有界连续仍不保证一致连续，例如 $\sin(1/x)$ 在 $(0,1)$ 上有界连续，却没有有限的左端点右极限。一致连续函数的和与差仍一致连续，乘积则需要额外控制；两个有界的一致连续函数的乘积由
\[
    |f(x)g(x)-f(y)g(y)|
    \leq\|f\|_\infty|g(x)-g(y)|+\|g\|_\infty|f(x)-f(y)|
\]
可知一致连续，其中 $\|f\|_\infty=\sup_{x\in I}|f(x)|$，$\|g\|_\infty$ 同理。不加有界条件时，$f(x)=g(x)=x$ 在 $\mathbb{R}$ 上均一致连续，但 $x^2$ 不一致连续：取 $x_n=n$、$y_n=n+1/n$，则两点距离趋于零，平方之差却趋于 $2$。

\section{无穷小量与无穷大量的阶}
\label{sec:orders-of-infinitesimals-and-infinities}

% 教材第 121--127 页；分章 PDF 第 39--45 页。
本节比较的对象是同一极限过程中变化的量，而不是某个固定的“小数”或“大数”。以下以 $x\to x_0$ 为规范；换成单侧趋近或趋于无穷远时，只需相应改变定义中的邻域条件。没有注明极限过程的阶比较是不完整的，例如 $x^2$ 与 $x$ 在 $x\to0$ 和 $x\to+\infty$ 时的相对大小恰好相反。

\begin{definition}
\label{def:infinitesimal-and-infinite-functions}
若 $f(x)\to0$，则称 $f$ 为该过程中的\emph{无穷小量}（infinitesimal）；若 $|f(x)|\to+\infty$，则称为\emph{无穷大量}（infinitely large quantity）。特别地，
\[
    f(x)\to A\in\mathbb{R}
    \quad\Longleftrightarrow\quad f(x)=A+\alpha(x),\quad\alpha(x)\to0.
\]
若 $f$ 在相应的去心邻域内恒不为零，则
\begin{equation}
    f(x)\to0
    \quad\Longleftrightarrow\quad
    \left|\frac1{f(x)}\right|\to+\infty.
    \label{eq:infinitesimal-and-reciprocal-infinity}
\end{equation}
\end{definition}

式~\eqref{eq:infinitesimal-and-reciprocal-infinity} 来自 $0<|f(x)|<1/M$ 与 $|1/f(x)|>M$ 的等价性，其中“恒不为零”保证倒数有定义。有限个无穷小量的和仍是无穷小量，无穷小量与局部有界函数的乘积也是无穷小量；这些均由极限运算或夹逼原理得到。无限多个无穷小量的相加则不能直接套用有限次运算的结论。

\begin{definition}
\label{def:little-o-big-o-and-asymptotic-equivalence}
设 $f,g$ 在相应去心邻域内有定义，且 $g$ 恒不为零。定义\emph{小 $o$ 记号}（little-o notation）、\emph{大 $O$ 记号}（big-O notation）与\emph{渐近等价}（asymptotic equivalence）如下：
\begin{equation}
    \begin{aligned}
        f(x)=o(g(x))
        &\quad\Longleftrightarrow\quad \frac{f(x)}{g(x)}\to0,\\
        f(x)=O(g(x))
        &\quad\Longleftrightarrow\quad
        \exists C>0,\ \exists\delta>0,\quad
        0<|x-x_0|<\delta\ \Longrightarrow\ |f(x)|\leq C|g(x)|,\\
        f(x)\sim g(x)
        &\quad\Longleftrightarrow\quad \frac{f(x)}{g(x)}\to1.
    \end{aligned}
    \label{eq:asymptotic-comparison-notations}
\end{equation}
若 $f,g$ 都是无穷小量且 $f/g\to0$，称 $f$ 是比 $g$ \emph{高阶的无穷小量}（higher-order infinitesimal）；若 $f/g\to c\in\mathbb{R}\setminus\{0\}$，称二者为\emph{同阶无穷小量}（infinitesimals of the same order），$c=1$ 时称为\emph{等价无穷小量}（equivalent infinitesimals）。若二者都是无穷大量，非零有限比值和比值为 $1$ 分别定义同阶与等价无穷大量，而 $f/g\to0$ 表示 $f$ 是比 $g$ \emph{低阶的无穷大量}。
\end{definition}

小 $o$ 与大 $O$ 本身不要求 $f,g$ 都趋于零或无穷。例如，$f=o(1)$ 就是 $f\to0$，$f=O(1)$ 表示局部有界。$f=o(g)$ 必推出 $f=O(g)$，反向不成立；按照本节采用的教材定义，同阶要求比值具有非零有限极限，比值只是在两个正常数之间波动并不足以称为同阶。例如，$f(x)=x(2+\sin(1/x))$ 与 $g(x)=x$ 在 $x\to0$ 时互为大 $O$，但比值没有极限。若 $f\sim g$，则 $f$ 也在充分小的去心邻域内不为零，且
\begin{equation}
    f\sim g
    \quad\Longleftrightarrow\quad
    f=g+o(g).
    \label{eq:asymptotic-equivalence-remainder}
\end{equation}

\emph[caution]{阶记号中的等号表达一种渐近关系，不能当作两个固定函数之间可任意倒转的恒等式。}同一行中不同的 $o(g)$ 可以代表不同函数。例如，$x^2=o(x)$ 也满足 $x^2=O(x)$，只能说明小 $o$ 类包含于大 $O$ 类，不能据此写成 $o(x)=O(x)$；$x=O(x)$ 却不满足 $x=o(x)$。当 $x\to0$ 时，正整数幂的指数越大，无穷小量阶越高；对倒数 $1/x^n$ 而言，指数越大则无穷大量阶越高。若需要非整数幂的双侧比较，使用 $|x-x_0|^\alpha$ 可避免负底数实幂的定义问题。

\begin{proposition}[对数、幂与指数的增长次序]
\label{prop:logarithm-power-exponential-orders}
设 $a>1$、$\alpha>0$。当 $x\to+\infty$ 时，
\begin{equation}
    \ln x=o(x^\alpha),
    \qquad x^\alpha=o(a^x).
    \label{eq:logarithm-power-exponential-orders}
\end{equation}
即对数增长慢于任意正幂，任意正幂增长慢于底数大于 $1$ 的指数函数。
\end{proposition}

\begin{proof}
先证第二式，取整数 $m>\alpha$，令 $n=\lfloor x\rfloor$、$h=a-1>0$。当 $n\geq2m$ 时，由二项式定理，
\[
    a^x\geq a^n=(1+h)^n\geq\binom nm h^m
    \geq\frac{h^m}{2^m m!}n^m.
\]
又 $x<n+1\leq2n$，所以
\[
    0<\frac{x^\alpha}{a^x}
    \leq\frac{2^{m+\alpha}m!}{h^m}n^{\alpha-m}\to0.
\]
对于第一式，令 $t=\ln x\to+\infty$，则
\[
    \frac{\ln x}{x^\alpha}=\frac{t}{(e^\alpha)^t}\to0,
\]
其中使用已经证明的幂与指数的比较，且 $e^\alpha>1$。
\end{proof}

\begin{theorem}[常用等价无穷小]
\label{thm:standard-equivalent-infinitesimals}
当 $x\to0$ 时，有
\begin{equation}
    \begin{aligned}
        \sin x&\sim x,&\tan x&\sim x,&1-\cos x&\sim\frac{x^2}{2},\\
        \arcsin x&\sim x,&\arctan x&\sim x,&\ln(1+x)&\sim x,\\
        e^x-1&\sim x,&a^x-1&\sim x\ln a,&(1+x)^\alpha-1&\sim\alpha x.
    \end{aligned}
    \label{eq:standard-equivalent-infinitesimals}
\end{equation}
其中 $a>0$、$a\neq1$，$\alpha\in\mathbb{R}\setminus\{0\}$；反三角函数取通常的主值分支，对数与实幂均在零点附近取其有定义的范围。
\end{theorem}

\begin{proof}
$\sin x\sim x$ 是第一个重要极限；由 $\cos x\to1$，得到 $\tan x/x=(\sin x/x)/\cos x\to1$。又
\[
    \frac{1-\cos x}{x^2/2}
    =\left(\frac{\sin(x/2)}{x/2}\right)^2\to1.
\]
令 $t=\arcsin x$，由反函数连续性知 $t\to0$，且 $x=\sin t$，故 $\arcsin x/x=t/\sin t\to1$；反正切的结论同理用 $x=\tan t$ 得到。

对第二个重要极限取连续的自然对数，得到
\[
    \frac{\ln(1+x)}x
    =\ln\bigl((1+x)^{1/x}\bigr)\to\ln e=1.
\]
令 $t=e^x-1$，则 $t\to0$ 且 $x=\ln(1+t)$，从而 $(e^x-1)/x=t/\ln(1+t)\to1$。再分别将 $x\ln a$ 和 $\alpha\ln(1+x)$ 代入指数等价式，得到
\[
    \frac{a^x-1}{x\ln a}\to1,
    \qquad
    \frac{(1+x)^\alpha-1}{\alpha x}
    =\frac{e^{\alpha\ln(1+x)}-1}{\alpha\ln(1+x)}
      \frac{\ln(1+x)}x\to1.
\]
参数非零保证所写比较量在去心邻域内不为零。
\end{proof}

结合式~\eqref{eq:asymptotic-equivalence-remainder}，这些等价式也可写成带余量的表达，例如
\[
    \sin x=x+o(x),\qquad
    \cos x=1-\frac{x^2}{2}+o(x^2),\qquad
    e^x=1+x+o(x),\qquad
    (1+x)^\alpha=1+\alpha x+o(x).
\]
最后一个展开在 $\alpha=0$ 时也成立，但此时不能把零函数写成与 $\alpha x$ 等价，因为比较的分母恒为零。所有这些式子都来自本章的极限与连续性，不需要提前使用泰勒公式。

\begin{theorem}[等价因子的替换]
\label{thm:replacement-of-equivalent-factors}
在同一极限过程中，设 $f\sim\widetilde f$、$g\sim\widetilde g$，所有函数在相应去心邻域内有定义，且商中的分母恒不为零。若右端极限存在，则
\begin{equation}
    \lim f(x)h(x)=\lim\widetilde f(x)h(x),
    \qquad
    \lim\frac{f(x)}{g(x)}=\lim\frac{\widetilde f(x)}{\widetilde g(x)}.
    \label{eq:replacement-of-equivalent-factors}
\end{equation}
这些结论对有限极限以及 $+\infty,-\infty$ 和绝对值趋于无穷的情形均成立。
\end{theorem}

\begin{proof}
由等价关系，$f/\widetilde f\to1$、$\widetilde g/g\to1$，因而
\[
    fh=\frac f{\widetilde f}\,(\widetilde f h),
    \qquad
    \frac fg=\frac f{\widetilde f}\frac{\widetilde g}g
        \frac{\widetilde f}{\widetilde g}.
\]
右端中额外的因子都趋于 $1$，有限极限情形由乘积法则成立。对无穷极限，这些因子最终为正，且可以限制在任意接近 $1$ 的正数区间内，所以也保持原来的无穷趋向。
\end{proof}

\begin{example}
\label{ex:equivalent-infinitesimal-calculations}
求 $\ln\cos2x$ 在 $x\to0$ 时形如 $cx^\alpha$ 的等价无穷小，并计算
\[
    \lim_{x\to0}\frac{\tan x-\sin x}{x^3}.
\]
\end{example}

\begin{solve}
在零点附近 $\cos2x>0$，且 $\cos2x=1-2\sin^2x$，故由 $\ln(1+t)\sim t$ 与 $\sin x\sim x$，
\[
    \ln\cos2x=\ln(1-2\sin^2x)\sim-2\sin^2x\sim-2x^2.
\]
第二问先作精确变形，使可能相消的差变成乘积：
\[
    \tan x-\sin x=\frac{\sin x(1-\cos x)}{\cos x}.
\]
再对乘积中的因子作等价替换，得到
\[
    \frac{\tan x-\sin x}{x^3}
    =\frac{\sin x}{x}\frac{1-\cos x}{x^2}\frac1{\cos x}
    \to1\cdot\frac12\cdot1=\frac12.
    \qedhere
\]
\end{solve}

\emph[caution]{等价替换可以用于乘除中的因子，不能无条件用于和差中的各项。}虽然 $\tan x\sim x$ 且 $\sin x\sim x$，但把上一例的分子直接替成 $x-x=0$ 就会丢失决定结果的三阶小量。等价只保证每项的相对误差趋于零；当主项相消时，这些误差可能与剩余量同阶甚至更大。处理和差应先恒等变形，或把余项写明并核对其相对于整个和差的阶。

\section*{习题三}
\addcontentsline{toc}{section}{习题三}
